\icode{N24L-A-SI-L2-260627-002}
\nceeID[3]{2024山东}{1}
2024年是中国航天大年，神舟十八号、嫦娥六号等已陆续飞天，部分航天器装载了具有抗干扰性强的核电池。已知\ce{^{90}_{38}Sr}衰变为\ce{^{90}_{39}Y}的半衰期约为29年；\ce{^{238}_{94}Pu}衰变为\ce{^{234}_{92}U}的半衰期约87年。现用相同数目的\ce{^{90}_{38}Sr}和\ce{^{238}_{94}Pu}各做一块核电池，下列说法正确的是\paren[]
\begin{choices}
	\item \ce{^{90}_{38}Sr}衰变为\ce{^{90}_{39}Y}时产生\(\alpha\)粒子
	\item \ce{^{238}_{94}Pu}衰变为\ce{^{234}_{92}U}时产生\(\beta\)粒子
	\item 50年后，剩余的\ce{^{90}_{38}Sr}数目大于\ce{^{238}_{94}Pu}的数目
	\item 87年后，剩余的\ce{^{90}_{38}Sr}数目小于\ce{^{238}_{94}Pu}的数目
\end{choices}
\begin{answer}
	D\quad 提示：\(N = N_0 (\frac{1}{2})^{\frac{t}{T}}\)
\end{answer}